\section{The square family and its affine group}
We use column coordinates and composition
$(A,t)(B,u)=(AB,t+Au)$. The source writes Euclidean transformations on
the right. Reversing an involutory word takes its inverse, so its order
and its defining power relation are unaffected by this convention.
Let $e_1,\ldots,e_4$ be the standard Euclidean basis, and put
\[
 h=\tfrac12(e_1+e_2+e_3+e_4),\qquad
 \Lambda=\Z e_1+\Z e_2+\Z e_3+\Z h.
\]
All integer vectors below are coordinates in the ordered basis
$(e_1,e_2,e_3,h)$ of $\Lambda$.

Take the five source reflections, here specified in Euclidean coordinates:
\begin{align*}
 R_0(x)&=(1-x_1,x_2,x_3,x_4),\\
 R_1(x)&=x-\tfrac12(x_1-x_2-x_3-x_4)(1,-1,-1,-1),\\
 R_2(x)&=(x_1,x_2,x_3,-x_4),\\
 R_3(x)&=(x_1,x_2,x_4,x_3),\\
 R_4(x)&=(x_1,x_3,x_2,x_4).
\end{align*}
Set $E=\langle R_0,\ldots,R_4\rangle$ and mark it by
\[
 t_0=R_1,\quad t_1=R_2,\quad t_2=R_1R_3,
 \quad t_3=R_0R_4.                                      \tag{S.1}
\]
These transformations preserve $\Lambda$. They generate $E$: besides
$R_1=t_0$ and $R_2=t_1$, one has
\[
 R_3=t_0t_2,\qquad R_4=(t_0t_3)^3,\qquad
 R_0=t_3(t_0t_3)^3.                                    \tag{S.2}
\]
The last two identities follow from the commuting pair $R_0,R_4$,
the commuting pair $R_1,R_4$ and $(R_1R_0)^3=1$.

For $J\subseteq\{0,1,2,3\}$ write $E_J=\langle t_j:j\in J\rangle$.
Let $Q_J$ be its linear image and $K_J$ its translation kernel.
The full linear image $Q$ has order $1152$. A faithful finite action
for its calculation is the action on the following $48$ Euclidean roots:
\[
 \{\pm e_i\}\ \cup\
 \{\tfrac12(\pm e_1\pm e_2\pm e_3\pm e_4)\}\ \cup\
 \{\pm e_i\pm e_j:i<j\}.                              \tag{S.3}
\]
The roots span the space, so fixing them all forces a linear map to
be the identity.

The exact affine computation gives the following complete table. The
two nonsingleton translation lattices have bases displayed below it.
\begin{center}\begin{tabular}{c|r|c@{\qquad}c|r|c}
\toprule
$J$&$|Q_J|$&$K_J$&$J$&$|Q_J|$&$K_J$\\\midrule
$\varnothing$&1&0&$\{3\}$&2&0\\
$\{0\}$&2&0&$\{0,3\}$&12&0\\
$\{1\}$&2&0&$\{1,3\}$&4&0\\
$\{0,1\}$&6&0&$\{0,1,3\}$&48&0\\
$\{2\}$&2&0&$\{2,3\}$&6&0\\
$\{0,2\}$&4&0&$\{0,2,3\}$&36&0\\
$\{1,2\}$&12&0&$\{1,2,3\}$&36&$K$\\
$\{0,1,2\}$&48&0&$\{0,1,2,3\}$&1152&$\Z^4$\\\bottomrule
\end{tabular}\end{center}
Here
\[
 K=\Z u+\Z v,\qquad
 u=(1,-1,1,0),\quad v=(0,-1,0,1).                    \tag{S.4}
\]
In particular the full translation lattice of $E$ is exactly $\Lambda$.
The entries are obtained by enumerating each finite $Q_J$, retaining
affine word lifts and computing all normalized Schreier translations,
as in AL.1 of the common paper. Integer Hermite reduction computes
their entire generated lattice. It is essential here to use normalized
transitions: a list of unnormalized differences alone would not suffice.

The table also exhibits an integral direct complement to $K$.
Writing $f_i$ for the integer coordinate vectors, the four columns
$u,v,f_2,f_3$ form a unimodular matrix. More explicitly, every
$x=(x_1,x_2,x_3,x_4)\in\Z^4$ has the unique expression
\[
 x=x_1u+x_4v+(x_2+x_1+x_4)f_2+(x_3-x_1)f_3.        \tag{S.5}
\]
Consequently $K\cap s\Z^4=sK$ for every positive integer $s$.

\section{The presentation and the period}
Put $a=t_0,b=t_1,c=t_2,d=t_3$. Let $\mathcal U$ be the abstract group
on four generators with their involution relations and
\begin{gather*}
 (ab)^3=(ac)^2=(ad)^6=(bc)^6=(bd)^2=(cd)^3=1,\\
 \bigl(c(da)^2\bigr)^2=
 \bigl(b(ad)^2a\bigr)^2=
 \bigl(a(bc)^2b\bigr)^2=1.                           \tag{S.6}
\end{gather*}
These are the square diagram relations and the three fixed toroidal
residue relations. All hold in the affine transformations. The marked
map $\mathcal U\to Q$ is onto. Its exact Reidemeister--Schreier
presentation, followed by Tietze simplification, has four kernel
generators and the following six relators:
\[
\begin{gathered}
 z_4^{-1}z_1^{-1}z_4z_1,\quad z_2z_1^{-1}z_2^{-1}z_1,
 \quad z_4^{-1}z_2z_4z_2^{-1},\\
 z_3z_1z_3^{-1}z_1^{-1},\quad z_3^{-1}z_4^{-1}z_3z_4,
 \quad z_3^{-1}z_2z_3z_2^{-1}.
\end{gathered}                                      \tag{S.7}
\]
Thus the actual kernel, not merely its abelianization, is free abelian
of rank four. The affine translation kernel has the same rank by the
preceding table. AL.2 now identifies $\mathcal U$ with $E$, with the
marking preserved.

Define $w=d(cb)^2$. Direct affine multiplication gives
\[
 w^2=(I,v).                                         \tag{S.8}
\]
The $Q$-orbit of $v$ consists of $24$ vectors and spans $\Z^4$ over
the integers. This is checked by its full integral Hermite basis $I_4$.
Conjugating a translation by an affine transformation applies its
linear part to the translation vector. Hence the normal closure of
$w^{2s}$ in $E$ is exactly $s\Lambda$, and
\[
 \mathcal U/\langle\!\langle w^{2s}\rangle\!\rangle
       \cong E/s\Lambda,
 \qquad |E/s\Lambda|=1152s^4.                       \tag{S.9}
\]
This proves a marked presentation identity for all $s$ at once.

For comparison with the source's original toroid, the five-generator
Coxeter presentation $[3,3,4,3]$ has the same affine realization.
The separate exact kernel calculation over $Q$ again reduces to four
generators with all six pairwise commutator relations and no others.
AL.2 identifies that presentation with $E$. Its period word satisfies
\[
 \bigl(R_0R_1(R_2R_3R_4)^3\bigr)^2=(I,f_4).
\]
The $Q$-orbit of $f_4$ also spans $\Z^4$. Its $s$th power therefore
imposes exactly $s\Lambda$. Together with (S.2), this proves that the
groups in Proposition 4.1 are the same marked constructions under
the displayed change of generators, not just groups of the same order.

\section{The ordered residues}
Let $P_i$ be the image of $E_{I\setminus\{i\}}$ in $E/s\Lambda$.
For $i=1,2,3$ the linear maps are injective on the entire affine
parabolics, so their groups do not collapse in any translation quotient.
In the respective ordered paths
\[
 (t_2,t_3,t_0),\qquad(t_1,t_0,t_3),\qquad(t_0,t_1,t_2),
\]
the edge orders are $3,6$ and the end generators commute. Their
additional relations in (S.6) are precisely those of
$\{3,6\}_{(1,1)},\{3,6\}_{(2,0)},\{3,6\}_{(2,0)}$.
The standard triangle presentations with these relations have orders
$36,48,48$, respectively; exact finite presentation calculations give
these orders. The table gives the same orders for the three images,
so the marked epimorphisms are isomorphisms. No classification by order
alone is being used.

For $P_0$, put $(A,B,C)=(t_3,t_2,t_1)$ and let
\[
 \Delta=\langle A,B,C\mid A^2=B^2=C^2=(AB)^3=(BC)^6=(AC)^2=1\rangle.
\]
Its epimorphism onto the actual linear group $Q_{\{1,2,3\}}$ of
order $36$ has kernel with the exact simplified presentation
\[
 \langle z_1,z_2\mid z_2z_1^{-1}z_2^{-1}z_1=1\rangle.
\]
The affine image has translation lattice $K$ of rank two. The
restriction from this kernel onto $K$ is a surjection of free abelian
groups of the same rank, hence an isomorphism. This is the same
rank argument as AL.2, now with rank two inside a four-dimensional
ambient affine space. Thus $E_{\{1,2,3\}}$ is the marked triangle group
$\Delta$ itself.

The $Q_{\{1,2,3\}}$-orbit of $v$ has six vectors and spans exactly $K$.
It follows that imposing $(A(BC)^2)^{2s}=1$ on $\Delta$ factors its
translation lattice by $sK$. By (S.5) the ambient quotient induces
exactly the same subgroup $sK$, without any extra residue collapse.
This is the standard marked presentation of $\{3,6\}_{(s,s)}$, and
its order is $36s^2$. The ordered rank-three groups are therefore
exactly the four claimed in Proposition 4.1.

The source also computes two word orders to distinguish the possible
residue periods. Here (S.8) has order $s$, while
$A(BC)^2B$ has linear part of order three and cube $(I,u)$.
The primitive vector $u$ has order $s$ modulo $s\Lambda$, so this
word has order $3s$. The word itself is generally not a translation:
it is its cube that is a translation. This corrects the terminology
in the source without changing either asserted order.

\section{Intersections and flag transitivity for every period}
The fixed linear marking has the intersection property. The certificate
checks all $16^2$ equalities $Q_J\cap Q_T=Q_{J\cap T}$ as literal
sets of matrices. Its four maximal parabolic coset sets have
$32,32,24,24$ members. Complete enumeration of pairwise incident cosets
gives the following numbers after translating the first coset to the
subgroup itself:
\[
\begin{array}{c|rrrrr}
\text{types}&012&013&023&123&0123\\\hline
\text{pairwise incident families}&18&18&18&18&36.
\end{array}
\]
Every listed family has a common element. Rank-one and rank-two
families already satisfy that assertion by nonemptiness and incidence.
This proves typed coset Helly for the fixed linear quotient. The
normalization loses no family, because right translation acts
transitively on the cosets of the first type.

Now fix a positive integer $s$ and set $G=E/s\Lambda$, with the
induced map $q:G\to Q$. Its kernel is $N=\Lambda/s\Lambda$.
For a proper marked parabolic $G_J$, the table and the fact that
$s\Lambda$ consists of translations give
\[
 N\cap G_J=
 \begin{cases}
 (K+s\Lambda)/s\Lambda,&J=\{1,2,3\},\\
 1,&\text{every other proper }J.
 \end{cases}                                       \tag{S.10}
\]
For the full index set the intersection is $N$. This entire kernel
profile has the intersection property: two different proper supports
cannot both have a nontrivial kernel intersection. UA.K therefore
proves the intersection property in $G$.

Equation (S.5) yields the internal direct product
\[
 N\cong (K/sK)\ \oplus\
       \bigl((\Z f_2+\Z f_3)/s(\Z f_2+\Z f_3)\bigr).
\]
Each maximal parabolic intersection in (S.10) is a product of some
of these two factors. The linear quotient has IP and typed coset
Helly, so UA.SK.Helly transfers typed coset Helly to $G$.
Its four markers have order exactly two, as their nontrivial linear
images do. UA.Geometry now proves that the coset incidence system
is thin, residually connected and flag-transitive, and that its full
type-preserving automorphism group is $G$.

Finally (S.6) and (S.9) give exactly the diagram relations and the four
ordered residue relations. There is no extra global relation. This
proves universality, completing Theorem 4.3 for its full stated range
$s\geq2$. The argument in fact also verifies these group and coset
assertions at $s=1$; no additional claim about the terminology for the
source's rank-five toroid at that parameter is needed.
